% Le birapport et le groupe anharmonique, d'après le dossier 147 de Grothendieck
% Compile avec pdflatex, xelatex ou lualatex.
\documentclass[11pt,a4paper]{article}

\usepackage{iftex}
\ifPDFTeX
  \usepackage[utf8]{inputenc}
  \usepackage[T1]{fontenc}
  \IfFileExists{lmodern.sty}{\usepackage{lmodern}}{}
  \usepackage[french]{babel}
\else
  \usepackage{fontspec}
  \usepackage{polyglossia}
  \setdefaultlanguage{french}
\fi
\usepackage[margin=2.6cm]{geometry}
\usepackage{amsmath,amssymb,amsthm,mathtools}
\usepackage{booktabs}
\usepackage{enumitem}
\usepackage{microtype}
\usepackage{tikz}
% Ce qu'une `tikzpicture` des fiches peut utiliser, en plus de TikZ.
%
% Lu deux fois : par chaque document (\input après \usepackage{tikz}), pour le
% PDF, et par scripts/tikz.mjs, qui compile chaque figure seule pour la vue web.
% Une bibliothèque ou une couleur ajoutée ici atteint les deux ; ajoutée
% ailleurs, elle n'en atteint qu'un, et le PDF et l'écran ne disent plus la
% même chose.

\usetikzlibrary{calc,arrows.meta,patterns,decorations.pathreplacing}

% La couleur d'accent du site, pour ce qu'une figure veut faire ressortir.
\definecolor{accent}{HTML}{38539D}
\definecolor{accent2}{HTML}{A04B2E}

\usepackage[hidelinks]{hyperref}
\emergencystretch=1.5em

% ---------- Environnements ----------
\theoremstyle{plain}
\newtheorem{theoreme}{Théorème}[section]
\newtheorem{proposition}[theoreme]{Proposition}
\newtheorem{lemme}[theoreme]{Lemme}
\newtheorem{corollaire}[theoreme]{Corollaire}
\theoremstyle{definition}
\newtheorem{definition}[theoreme]{Définition}
\newtheorem{exemple}[theoreme]{Exemple}
\newtheorem{exercice}{Exercice}
\theoremstyle{remark}
\newtheorem{remarque}[theoreme]{Remarque}

% ---------- Notations ----------
\newcommand{\C}{\mathbb{C}}
\newcommand{\Z}{\mathbb{Z}}
\newcommand{\Pu}{\mathbb{P}^1}
\newcommand{\PGL}{\mathrm{PGL}}
\newcommand{\Sym}{\mathfrak{S}}

\title{Le birapport et le groupe anharmonique\\[2pt]
\large Quatre points sur une droite, et ce qu'en garde l'action de \texorpdfstring{$\Sym_4$}{S4}}
\author{Michel Hua, d'après le dossier 147 des archives de Grothendieck}
\date{Octobre 2026}

\begin{document}
\maketitle

\begin{abstract}
Quatre points distincts de la droite projective complexe ont un invariant, leur birapport, qui ne change pas par homographie et qui les caractérise à homographie près. Quand on permute les quatre points, le birapport prend six valeurs au plus : le groupe $\Sym_4$ agit à travers un quotient isomorphe à $\Sym_3$, le groupe anharmonique, qui opère sur $\Pu\setminus\{0,1,\infty\}$ par $\lambda\mapsto1-\lambda$, $\lambda\mapsto1/\lambda$, etc. Cette action a exactement deux orbites exceptionnelles, celle des quadruplets harmoniques et celle des équianharmoniques, et ses orbites sont les fibres d'une seule fraction rationnelle, l'invariant $j$. Grothendieck s'en sert, dans des notes du début des années 1980 sur la « structure à l'infini » des espaces de modules de courbes, pour décrire l'espace des sphères à quatre points marqués et calculer sa caractéristique d'Euler. On démontre ici tout ce que ce calcul utilise, avec la théorie élémentaire des actions de groupe.
\end{abstract}

\paragraph{Leçons concernées.} 101 (Groupe opérant sur un ensemble. Exemples et applications), et en partie 182 (Applications des nombres complexes à la géométrie) et 191 (Exemples d'utilisation de techniques d'algèbre en géométrie).

\paragraph{Développements possibles.} L'action de $\Sym_4$ sur les birapports, son noyau $V_4$ et le groupe anharmonique (théorème \ref{thm:anharmonique}) ; les orbites exceptionnelles et l'invariant $j$ (théorème \ref{thm:j}). Le fichier \texttt{BirapportAnharmonique.lean} vérifie avec Mathlib l'invariance du birapport par les homographies (proposition \ref{prop:invariance}), la table du groupe anharmonique, les points fixes (proposition \ref{prop:orbites}), l'invariance de $j$ et sa factorisation (théorème \ref{thm:j}), et le calcul de la caractéristique d'Euler de la remarque \ref{rem:euler}.

\paragraph{Prérequis.} Homographies et droite projective $\Pu(\C)=\C\cup\{\infty\}$, actions de groupe et formule orbite-stabilisateur, fractions rationnelles.

\paragraph{Notations.} $\PGL_2(\C)$ agit sur $\Pu(\C)$ par homographies $z\mapsto\frac{az+b}{cz+d}$. On note $j_3=e^{2i\pi/3}$, et $V_4=\{\mathrm{id},(12)(34),(13)(24),(14)(23)\}$ le groupe de Klein, distingué dans $\Sym_4$.

\tableofcontents

% ==========================================================
\section{Le birapport}
% ==========================================================

Une homographie est déterminée par l'image de trois points : le groupe $\PGL_2(\C)$ agit de façon simplement transitive sur les triplets de points distincts de $\Pu(\C)$. Pour comparer deux quadruplets, on envoie donc les trois premiers points sur une position de référence, $\infty$, $0$, $1$, et l'on regarde où va le quatrième.

\begin{definition}
Le \emph{birapport} de quatre points distincts $z_1,z_2,z_3,z_4$ de $\Pu(\C)$ est l'image de $z_4$ par l'unique homographie $h$ telle que $h(z_1)=\infty$, $h(z_2)=0$, $h(z_3)=1$. Quand les $z_i$ sont finis,
\[
  [z_1,z_2;z_3,z_4]=\frac{(z_3-z_1)(z_4-z_2)}{(z_3-z_2)(z_4-z_1)} ,
\]
et l'on prolonge la formule par continuité quand l'un des points est $\infty$.
\end{definition}

En effet $h(z)=\frac{(z-z_2)(z_3-z_1)}{(z-z_1)(z_3-z_2)}$ envoie $z_1,z_2,z_3$ sur $\infty,0,1$. Comme les points sont distincts, le birapport est différent de $0$, $1$ et $\infty$.

\begin{proposition}\label{prop:invariance}
Le birapport est invariant par homographie : $[g(z_1),g(z_2);g(z_3),g(z_4)]=[z_1,z_2;z_3,z_4]$ pour tout $g\in\PGL_2(\C)$. Deux quadruplets de points distincts sont dans la même orbite de $\PGL_2(\C)$ si et seulement s'ils ont le même birapport.
\end{proposition}

\begin{proof}
Si $h$ envoie $z_1,z_2,z_3$ sur $\infty,0,1$, alors $h\circ g^{-1}$ envoie $g(z_1),g(z_2),g(z_3)$ sur $\infty,0,1$, et $g(z_4)$ sur $h(z_4)$ : c'est l'invariance. Réciproquement, si deux quadruplets ont le même birapport $\lambda$, les homographies $h$ et $h'$ correspondantes les envoient tous deux sur $(\infty,0,1,\lambda)$, et $h'^{-1}\circ h$ envoie le premier sur le second.
\end{proof}

Ainsi l'application « birapport » identifie l'ensemble des quadruplets ordonnés de points distincts, modulo homographie, à $\Pu(\C)\setminus\{0,1,\infty\}$. C'est l'égalité que Grothendieck écrit $\mathcal M_{0,4}\simeq\Pu\setminus\{0,1,\infty\}$ : l'espace des sphères munies de quatre points numérotés est la sphère privée de trois points.

\begin{exemple}\label{ex:configurations}
Les sommets $1,i,-1,-i$ d'un carré ont pour birapport $[1,i;-1,-i]=2$. Le centre et les sommets d'un triangle équilatéral, $0,1,j_3,j_3^2$, ont pour birapport $[0,1;j_3,j_3^2]=e^{-i\pi/3}$. On verra que ces deux nombres sont exceptionnels (figure \ref{fig:configurations}).
\end{exemple}

\begin{figure}[h]
\centering
\begin{tikzpicture}[scale=1.3,pt/.style={circle,fill,inner sep=1.4pt}]
  \begin{scope}
    \draw[gray!50] (0,0) circle (1);
    \draw[accent,thick] (1,0) -- (0,1) -- (-1,0) -- (0,-1) -- cycle;
    \foreach \p/\l/\pos in {{(1,0)}/z_1=1/right,{(0,1)}/z_2=i/above,{(-1,0)}/z_3=-1/left,{(0,-1)}/z_4=-i/below} \node[pt,label={\pos:$\l$}] at \p {};
    \node[font=\small] at (0,-1.75) {harmonique : $\lambda=2$};
  \end{scope}
  \begin{scope}[xshift=4.4cm]
    \draw[gray!50] (0,0) circle (1);
    \draw[accent2,thick] (1,0) -- (-0.5,0.866) -- (-0.5,-0.866) -- cycle;
    \draw[accent2,dashed] (0,0) -- (1,0) (0,0) -- (-0.5,0.866) (0,0) -- (-0.5,-0.866);
    \foreach \p/\l/\pos in {{(0,0)}/z_1=0/below right,{(1,0)}/z_2=1/right,{(-0.5,0.866)}/z_3=j_3/above left,{(-0.5,-0.866)}/z_4=j_3^2/below left} \node[pt,label={\pos:$\l$}] at \p {};
    \node[font=\small] at (0,-1.75) {équianharmonique : $\lambda=e^{-i\pi/3}$};
  \end{scope}
\end{tikzpicture}
\caption{Les deux configurations exceptionnelles de quatre points : un carré, et un triangle équilatéral avec son centre. Leurs symétries expliquent leurs stabilisateurs, d'ordre $2$ et $3$ dans le groupe anharmonique.}\label{fig:configurations}
\end{figure}

% ==========================================================
\section{Permuter les quatre points}
% ==========================================================

Le groupe $\Sym_4$ agit sur les quadruplets ordonnés en permutant les points, et cette action commute avec celle des homographies. Elle passe donc au quotient : $\Sym_4$ agit sur $\Pu\setminus\{0,1,\infty\}$, par $\pi\cdot[z_1,z_2;z_3,z_4]=[z_{\pi^{-1}(1)},\dots,z_{\pi^{-1}(4)}]$. Le calcul direct donne six valeurs.

\begin{lemme}\label{lem:six}
Si $\lambda=[z_1,z_2;z_3,z_4]$, alors
\[
  [z_2,z_1;z_3,z_4]=\frac1\lambda,\qquad [z_1,z_3;z_2,z_4]=1-\lambda,
\]
et les permutations de $z_1,z_2,z_3,z_4$ donnent les six valeurs
\[
  \lambda,\quad \frac1\lambda,\quad 1-\lambda,\quad \frac1{1-\lambda},\quad \frac{\lambda-1}\lambda,\quad \frac\lambda{\lambda-1} .
\]
Une double transposition ne change pas le birapport.
\end{lemme}

\begin{proof}
Échanger $z_1$ et $z_2$ inverse la formule. Échanger $z_2$ et $z_3$ donne $\frac{(z_2-z_1)(z_4-z_3)}{(z_2-z_3)(z_4-z_1)}$, et l'identité $(z_3-z_1)(z_4-z_2)-(z_2-z_1)(z_4-z_3)=(z_3-z_2)(z_4-z_1)$, qui se vérifie en développant, montre que la somme des deux birapports vaut $1$. Pour $(12)(34)$, la formule devient $\frac{(z_4-z_2)(z_3-z_1)}{(z_4-z_1)(z_3-z_2)}$, qui est la même ; de même pour $(13)(24)$, et $(14)(23)$ est leur produit. Les transpositions $(12)$ et $(23)$ engendrent $\Sym_4$ modulo $V_4$, et leurs actions $\lambda\mapsto1/\lambda$ et $\lambda\mapsto1-\lambda$ engendrent les six fonctions de l'énoncé.
\end{proof}

\begin{theoreme}[le groupe anharmonique]\label{thm:anharmonique}
Les six fonctions du lemme \ref{lem:six} forment, pour la composition, un groupe $A$ isomorphe à $\Sym_3$. L'action de $\Sym_4$ sur les birapports se fait à travers un morphisme surjectif $\Sym_4\to A$ de noyau $V_4$, d'où $\Sym_4/V_4\simeq\Sym_3$.
\end{theoreme}

\begin{proof}
Chaque fonction du lemme est une homographie qui permute $\{0,1,\infty\}$ : par exemple $\lambda\mapsto1-\lambda$ échange $0$ et $1$ et fixe $\infty$, et $\lambda\mapsto\frac1{1-\lambda}$ envoie $0\mapsto1\mapsto\infty\mapsto0$. Une homographie qui fixe $0$, $1$ et $\infty$ est l'identité, donc l'application qui associe à une fonction sa permutation de $\{0,1,\infty\}$ est un morphisme injectif de $A$ dans $\Sym_3$ ; il est bijectif, car $A$ a six éléments, et $A$ est bien un groupe, image réciproque d'un groupe. L'action de $\Sym_4$ fournit un morphisme $\Sym_4\to A$, surjectif d'après le lemme, et dont le noyau contient $V_4$ ; par cardinalité ($24=4\cdot6$), le noyau est exactement $V_4$.
\end{proof}

\begin{remarque}
Le morphisme $\Sym_4\to\Sym_3$ a une description sans birapport : $\Sym_4$ permute les trois façons de partager $\{1,2,3,4\}$ en deux paires, $\{12|34\}$, $\{13|24\}$, $\{14|23\}$, et le noyau de cette action est $V_4$. Grothendieck lit ces trois partitions comme les trois « points à l'infini » de $\overline{\mathcal M}_{0,4}$, les valeurs $0$, $1$, $\infty$ du birapport, où la sphère se casse en deux sphères portant chacune deux des quatre points.
\end{remarque}

% ==========================================================
\section{Les orbites}
% ==========================================================

Le groupe anharmonique $A\simeq\Sym_3$ agit sur $U=\Pu(\C)\setminus\{0,1,\infty\}$. La formule orbite-stabilisateur dit qu'une orbite a $6$, $3$, $2$ ou $1$ éléments selon que le stabilisateur est trivial, d'ordre $2$, d'ordre $3$ ou $A$ tout entier.

\begin{proposition}\label{prop:orbites}
L'action de $A$ sur $U$ a exactement deux orbites non libres :
\begin{enumerate}[nosep]
\item l'orbite \emph{harmonique} $\{-1,\frac12,2\}$, où le stabilisateur est d'ordre $2$ ;
\item l'orbite \emph{équianharmonique} $\{e^{i\pi/3},e^{-i\pi/3}\}=\{-j_3^2,-j_3\}$, où le stabilisateur est d'ordre $3$.
\end{enumerate}
Toutes les autres orbites ont six éléments.
\end{proposition}

\begin{proof}
On cherche les points fixes des cinq éléments non triviaux de $A$, dans $U$. Pour les trois involutions : $1-\lambda=\lambda$ donne $\lambda=\frac12$ ; $\frac1\lambda=\lambda$ donne $\lambda=\pm1$, donc $\lambda=-1$ ; $\frac\lambda{\lambda-1}=\lambda$ donne $\lambda^2=2\lambda$, donc $\lambda=2$. Pour les deux éléments d'ordre $3$, $\frac1{1-\lambda}=\lambda$ et $\frac{\lambda-1}\lambda=\lambda$ équivalent tous deux à $\lambda^2-\lambda+1=0$, dont les racines sont $e^{\pm i\pi/3}$. Ces cinq points sont les seuls à avoir un stabilisateur non trivial, et ils se répartissent en une orbite de trois éléments ($\lambda\mapsto1-\lambda$ envoie $-1$ sur $2$, $\lambda\mapsto1/\lambda$ envoie $2$ sur $\frac12$) et une orbite de deux ($\lambda\mapsto1/\lambda$ échange $e^{i\pi/3}$ et $e^{-i\pi/3}$).
\end{proof}

Les quadruplets de l'exemple \ref{ex:configurations} sont dans ces orbites : un carré est harmonique, un triangle équilatéral avec son centre est équianharmonique. Le stabilisateur reflète leurs symétries. L'image réciproque dans $\Sym_4$ d'un stabilisateur d'ordre $2$ a $2\cdot4=8$ éléments : c'est le groupe diédral des symétries du carré. Celle d'un stabilisateur d'ordre $3$ a $12$ éléments : c'est le groupe alterné $\mathfrak A_4$, groupe des rotations d'un tétraèdre régulier, dont les sommets, sur la sphère de Riemann, sont justement un triangle équilatéral et son centre vus depuis le quatrième sommet.

\begin{figure}[h]
\centering
\begin{tikzpicture}[scale=1.9,pt/.style={circle,fill,inner sep=1.3pt}]
  \draw[->,gray] (-1.7,0) -- (2.7,0);
  \draw[->,gray] (0,-1.25) -- (0,1.35);
  \draw[gray!60] (0,0) circle (1);
  \draw[gray!60] (1,0) circle (1);
  \draw[gray!60] (0.5,-1.25) -- (0.5,1.35);
  \node[pt,fill=white,draw] at (0,0) {};
  \node[pt,fill=white,draw] at (1,0) {};
  \node[below left,font=\small] at (0,0) {$0$};
  \node[below right,font=\small] at (1,0) {$1$};
  \foreach \x/\l in {-1/-1,0.5/\frac12,2/2} \node[pt,accent] at (\x,0) {};
  \node[accent,below left,font=\small] at (-1,0) {$-1$};
  \node[accent,below right,font=\small] at (0.5,0) {$\frac12$};
  \node[accent,below right,font=\small] at (2,0) {$2$};
  \node[pt,accent2] at (0.5,0.866) {};
  \node[pt,accent2] at (0.5,-0.866) {};
  \node[accent2,above right,font=\small] at (0.5,0.866) {$e^{i\pi/3}$};
  \node[accent2,below right,font=\small] at (0.5,-0.866) {$e^{-i\pi/3}$};
\end{tikzpicture}
\caption{Les points exceptionnels de l'action anharmonique : l'orbite harmonique (en bleu) et l'orbite équianharmonique (en brun), aux intersections de la droite $\operatorname{Re}\lambda=\frac12$ et des cercles $|\lambda|=1$, $|\lambda-1|=1$, que le groupe anharmonique permute. Ces trois courbes découpent $\C\setminus\{0,1\}$ en domaines échangés par l'action.}\label{fig:orbites}
\end{figure}

% ==========================================================
\section{L'invariant \texorpdfstring{$j$}{j}}
% ==========================================================

Pour séparer les orbites, il suffit d'une fonction invariante qui prenne des valeurs distinctes sur des orbites distinctes. On la fabrique à partir des fonctions symétriques des six valeurs du lemme \ref{lem:six}.

\begin{theoreme}\label{thm:j}
La fraction rationnelle
\[
  j(\lambda)=256\,\frac{(\lambda^2-\lambda+1)^3}{\lambda^2(\lambda-1)^2}
\]
est invariante par le groupe anharmonique, et deux points de $U$ ont même image par $j$ si et seulement s'ils sont dans la même orbite. De plus
\[
  j(\lambda)-1728=64\,\frac{(\lambda+1)^2(\lambda-2)^2(2\lambda-1)^2}{\lambda^2(\lambda-1)^2},
\]
de sorte que $j$ vaut $1728$ sur l'orbite harmonique et $0$ sur l'orbite équianharmonique.
\end{theoreme}

\begin{proof}
L'invariance par $\lambda\mapsto1-\lambda$ est visible ($\lambda^2-\lambda+1$ et $\lambda(\lambda-1)$ sont invariants). Pour $\lambda\mapsto1/\lambda$, on multiplie numérateur et dénominateur par $\lambda^6$ : $(\lambda^{-2}-\lambda^{-1}+1)^3\lambda^6=(1-\lambda+\lambda^2)^3$ et $\lambda^{-2}(\lambda^{-1}-1)^2\lambda^6=(1-\lambda)^2\lambda^2$. Ces deux transformations engendrent $A$.

Soit $c\in\C$. L'équation $j(\lambda)=c$ s'écrit $256(\lambda^2-\lambda+1)^3-c\,\lambda^2(\lambda-1)^2=0$, équation polynomiale de degré $6$ en $\lambda$, dont aucune racine n'est $0$ ou $1$ (en ces points, le premier terme vaut $256$). Elle a donc au plus six solutions dans $U$, et elle contient l'orbite de chacune d'elles. Si l'une des solutions a une orbite à six éléments, les solutions forment exactement cette orbite. Sinon, toutes les solutions ont un stabilisateur non trivial, donc sont parmi les cinq points de la proposition \ref{prop:orbites}, et l'identité de l'énoncé, qui se vérifie en développant, montre que $j$ vaut $1728$ sur l'orbite harmonique (racines doubles) et $0$ sur l'orbite équianharmonique (racines triples) : les deux orbites exceptionnelles ont des images distinctes. Dans tous les cas, la fibre de $c$ est une seule orbite.
\end{proof}

\begin{remarque}[courbes elliptiques]
La courbe d'équation $y^2=x(x-1)(x-\lambda)$, sous la forme de Legendre, est une courbe elliptique dont les quatre points de $2$-torsion se projettent sur $0$, $1$, $\lambda$ et $\infty$. Changer l'ordre de ces quatre points change $\lambda$ par le groupe anharmonique, et $j(\lambda)$ est l'invariant $j$ de la courbe, qui la caractérise à isomorphisme près. Les orbites exceptionnelles correspondent aux courbes qui ont plus d'automorphismes que $\pm1$ : $j=1728$ pour la courbe harmonique $y^2=x^3-x$, d'automorphismes $\Z/4$, et $j=0$ pour la courbe équianharmonique $y^2=x^3-1$, d'automorphismes $\Z/6$. C'est le lien que Grothendieck fait, dans les mêmes pages, entre les sphères à quatre points et les courbes elliptiques.
\end{remarque}

\begin{remarque}[une caractéristique d'Euler]\label{rem:euler}
Grothendieck calcule la caractéristique d'Euler de l'espace des quatre points non numérotés en comptant chaque point avec le poids $1/|\text{stabilisateur}|$. La sphère privée de trois points a pour caractéristique $\chi(U)=2-3=-1$. On la découpe en parties stables par $A$ : les deux points équianharmoniques, de stabilisateur d'ordre $3$, les trois points harmoniques, de stabilisateur d'ordre $2$, et le reste, de caractéristique $-1-2-3=-6$, où $A$ agit librement. Le poids total vaut
\[
  \frac{-6}6+\frac{3}{6}+\frac26=-1+\frac12+\frac13=-\frac16=\frac{\chi(U)}{|A|},
\]
ce que la formule orbite-stabilisateur garantit : sur une orbite $O$, $\sum_{x\in O}\frac1{|A|}=\frac1{|\mathrm{Stab}(x)|}$. Comme $\Sym_4$ agit à travers $A$ avec un noyau $V_4$ d'ordre $4$ qui fixe tout, il faut encore diviser par $4$ pour obtenir la caractéristique relative à $\Sym_4$ : $-\frac1{24}$. La page écrit « $-1+\frac13+\frac12=-\frac1{24}$ », ce qui ne fait pas le compte : la somme vaut $-\frac16$, et le facteur $\frac14$ du noyau $V_4$ est omis. La valeur finale, que la marge donne aussi, est juste.
\end{remarque}

% ==========================================================
\section{Exercices}
% ==========================================================

\begin{exercice}\label{exo:carre}
Vérifier les deux birapports de l'exemple \ref{ex:configurations}. Plus généralement, montrer que quatre points d'un cercle sont harmoniques si et seulement si, dans un ordre convenable, ils forment un quadrilatère dont les produits des côtés opposés sont égaux.
\end{exercice}

\begin{exercice}\label{exo:reel}
Montrer que le birapport de quatre points distincts est réel si et seulement si les quatre points sont sur un même cercle ou une même droite.
\end{exercice}

\begin{exercice}\label{exo:car3}
On remplace $\C$ par un corps $K$ de caractéristique $3$. Montrer que les deux orbites exceptionnelles se confondent : $-1$ est alors racine double de $\lambda^2-\lambda+1$, et $j(-1)=0=1728$. Que se passe-t-il en caractéristique $2$ ?
\end{exercice}

\begin{exercice}\label{exo:partitions}
Montrer directement que l'action de $\Sym_4$ sur les trois partitions de $\{1,2,3,4\}$ en deux paires définit un morphisme surjectif $\Sym_4\to\Sym_3$ de noyau $V_4$.
\end{exercice}

% ==========================================================
\section{Corrigés}
% ==========================================================

\paragraph{Exercice 1.} $[1,i;-1,-i]=\frac{(-1-1)(-i-i)}{(-1-i)(-i-1)}=\frac{4i}{(1+i)^2}=\frac{4i}{2i}=2$. Et $[0,1;j_3,j_3^2]=\frac{j_3(j_3^2-1)}{(j_3-1)j_3^2}=\frac{j_3+1}{j_3}=1+j_3^2=-j_3=e^{-i\pi/3}$, en utilisant $j_3^2-1=(j_3-1)(j_3+1)$ et $1+j_3+j_3^2=0$. Pour quatre points d'un cercle, le module du birapport est $\frac{|z_3-z_1||z_4-z_2|}{|z_3-z_2||z_4-z_1|}$, et le birapport est réel (exercice 2). Il vaut $-1$, ou une autre valeur de l'orbite harmonique selon l'ordre, si et seulement si son module vaut $1$ dans l'ordre où il est négatif, c'est-à-dire si les produits de côtés opposés sont égaux.

\paragraph{Exercice 2.} Par une homographie, qui préserve à la fois le birapport et l'ensemble des cercles et droites de $\Pu(\C)$, on se ramène à $z_1=\infty$, $z_2=0$, $z_3=1$. Le birapport vaut alors $z_4$, et il est réel si et seulement si $z_4$ est sur la droite réelle, qui est le cercle généralisé passant par $\infty$, $0$ et $1$.

\paragraph{Exercice 3.} En caractéristique $3$, $\lambda^2-\lambda+1=\lambda^2+2\lambda+1=(\lambda+1)^2$, et les trois points $-1$, $\frac12$ et $2$ sont tous égaux à $2=-1$ (car $2\cdot2=4=1$) : l'orbite harmonique devient le seul point $\lambda=-1$, qui est aussi la racine double équianharmonique. Le numérateur de $j$ s'annule en $-1$, et $1728=3^3\cdot64=0$. En caractéristique $2$, $\frac12$ n'a pas de sens, $-1=1$ est exclu, et l'orbite harmonique disparaît ; c'est pourquoi Grothendieck limite ses formules à $\mathbb Z[\frac16]$, en inversant $2$ et $3$.

\paragraph{Exercice 4.} Une permutation envoie une partition en deux paires sur une partition en deux paires, d'où une action, donc un morphisme $\Sym_4\to\Sym_3$. Une double transposition, par exemple $(12)(34)$, fixe chaque partition : $\{12|34\}$, $\{13|24\}\mapsto\{24|13\}$, $\{14|23\}\mapsto\{23|14\}$. La transposition $(12)$ fixe $\{12|34\}$ et échange les deux autres, et le $3$-cycle $(234)$ les permute circulairement : l'image contient une transposition et un $3$-cycle, c'est $\Sym_3$. Le noyau a donc $24/6=4$ éléments et contient $V_4$ : c'est $V_4$.

\begin{thebibliography}{9}
\bibitem{g147} A. Grothendieck, \emph{Structure à l'infini des $M_{g,\nu}$}, notes manuscrites (début des années 1980), dossier 147 des archives de Grothendieck (université de Montpellier) ; transcription et lecture modernisée sur \url{https://grothendieck.commutator.io}.
\bibitem{audin} M. Audin, \emph{Géométrie}, EDP Sciences, 2006, chapitre VI (géométrie projective, birapport).
\bibitem{berger} M. Berger, \emph{Géométrie}, Nathan, 1990, chapitre 6.
\bibitem{silverman} J. H. Silverman, \emph{The Arithmetic of Elliptic Curves}, Springer, 2009, chapitre III, § 1 (forme de Legendre et invariant $j$).
\end{thebibliography}

\paragraph{Source.} Le dossier 147 des archives de Grothendieck, notes sur la structure à l'infini des espaces de modules de courbes, lu dans la transcription et la lecture modernisée du dépôt grothendieck-project, qui n'ont pas encore été relues sur les pages par un humain. La fiche démontre les faits élémentaires que le dossier utilise pour la sphère à quatre points, et corrige une somme de son calcul de caractéristique d'Euler.

\end{document}
